% Current implementation and probability ordering are source-audited for v5.0.1.
\subsection{A GP calculation of scan-global significance}
\label{v5-sec:global-method}\label{sec:v5-global}

A local probability describes a fluctuation at a fixed mass. A global
probability asks how often a complete background experiment would give
an equally extreme result anywhere in the declared search. Neighboring
hypotheses reuse events and sidebands, so a fluctuation can affect several
fits. The method of Ananiev and Read~\cite{AnanievRead2023} estimates the
correlation of the local-significance field from controlled background
perturbations, then generates many inexpensive correlated fields.

\methodfig{0.99}{../figures/v501_global_gp_flow.pdf}{The computational gain in the
significance-GP method. Full analysis scans of controlled perturbations
estimate the field covariance once. Gaussian vector draws then produce
many correlated scans without repeating those fits. A separate ensemble
of complete Poisson experiments tests the approximation. This is an
explanatory diagram of the archived HPS implementation, following
Ref.~\cite{AnanievRead2023}; the offset and width are retained explicitly.}
{fig:v501-global-gp-flow}

The background-interpolation GP predicts counts inside a mass window.
The significance GP instead approximates how the \emph{fit result varies
across tested masses} in repeated background experiments. Let $\psi$ be the
single-campaign yield or common coupling and let $\ell_p$ be its profiled
negative log likelihood. The signed local coordinate is
\begin{equation}
 r(\bm n;m)=\operatorname{sgn}(\widehat\psi^{\rm unc})
 \sqrt{2\bigl[\ell_p(0;m)-\ell_p(\widehat\psi^{\rm unc};m)\bigr]}.
 \label{v5-eq:signed-local-mapping}
\end{equation}
It retains both excesses and deficits. The conventional one-sided
asymptotic display uses
\begin{equation}
 Z_{\rm local}=\max(0,r),\qquad p_0=1-\Phi(Z_{\rm local}).
\end{equation}
A negative $r$ therefore has zero excess significance, with $p_0=0.5$
under this display convention.

\paragraph{From the analysis response to a covariance.}
For one coherent expected spectrum $\bm B$, first evaluate the unfluctuated
scan. Then increase one support bin by its Poisson standard deviation and
repeat the complete scan:
\begin{align}
 a_m&=r(\bm B;m),\\
 D_{im}&=r(\bm B+\sqrt{B_i}\,\bm e_i;m)-a_m.
 \label{eq:v501-response-directions}
\end{align}
Here $\bm e_i$ changes only bin $i$. The offset $a_m$ is the response to
the deterministic reference spectrum; it need not equal the exact mean
of the Poisson-fit distribution. Each row of $D$ represents one
independent noise direction, not one random Monte Carlo experiment.
The covariance, marginal width and correlation are
\begin{align}
 \Gamma_{mn}&=\sum_i D_{im}D_{in},
 \label{v5-eq:field-response}\\
 s_m&=\sqrt{\Gamma_{mm}},\\
 R_{mn}&=\frac{\Gamma_{mn}}{s_ms_n}.
 \label{eq:v501-field-correlation}
\end{align}
There is no division by the number of support bins. To first order,
$D_{im}\simeq\sqrt{B_i}\,\partial r_m/\partial n_i$, so the sum propagates
independent Poisson-bin noise through the full fitting response. The
matrix $R$ is estimated from that response; no RBF kernel is imposed on
the significance field.

\paragraph{Correlated simulations and one declared ordering.}
A simulated field has
\begin{equation}
 \bm z^*\sim\mathcal N(\bm0,R),\qquad
 r_m^*=a_m+s_m z_m^*.
 \label{eq:v501-field-draw}
\end{equation}
The archived extension retains a nonzero offset and nonunit width rather
than assuming an exactly centered, unit-normal marginal distribution.
Subtracting $a_m$ when constructing the response centers the covariance;
it does not subtract events from the observed spectrum.

The principal ordering standardizes each coordinate and keeps only
positive raw signal fits:
\begin{align}
 z_m&=\frac{r_m-a_m}{s_m},\\
 t_m&=\begin{cases}z_m,&r_m>0,\\-\infty,&r_m\leq0,\end{cases}
 &p_m^{\rm cond}&=\begin{cases}1-\Phi(z_m),&r_m>0,\\1,&r_m\leq0.\end{cases}
 \label{eq:v501-local-ordering}
\end{align}
For the fixed full search grid $\mathcal M$, define
\begin{equation}
 T=\max_{m\in\mathcal M}t_m,
 \qquad
 p_{\rm global}=\Pr(T^*\geq T_{\rm obs}).
 \label{v5-eq:global-ordering}
\end{equation}
If no fitted signal is positive, $T=-\infty$. For $N$ complete simulated
fields, the estimated tail is
\begin{equation}
 \widehat p_{\rm global}=\frac{k}{N},
 \qquad k=\sum_{j=1}^{N}\mathbf1[T_j^*\geq T_{\rm obs}].
 \label{eq:v501-global-mc-tail}
\end{equation}
At zero exceedances the one-sided 95\% sampling bound is
$1-0.05^{1/N}$; no zero probability is assigned. A pointwise global curve
uses each observed $t_m$ as the threshold against the \emph{same full-grid}
maximum distribution.

A separate archived ordering uses $U=\max_m\max(0,r_m)$. It tests the
largest raw fluctuation, whereas $T$ ranks positive fits relative to the
chosen background's offset and width. These probabilities cannot be
interchanged after inspecting their values. The positive-fit gate also
explains why the conditional convention gives one for nonpositive fits,
while the nominal asymptotic plot gives 0.5.

Each individual study began with ten complete Poisson scans, followed by
1,000 independently generated complete scans and 200,000 GP fields.
Every mass in one scan uses the same full spectrum. Sideband targets,
count-dependent errors and posterior predictions are updated; the reviewed
kernel coordinates remain fixed. Independently generated pointwise toy
collections cannot be joined by their row numbers to reconstruct this
correlation. The 2015, 2016 and 2021 grids contain 72, 142 and 201 integer
mass hypotheses over 19--90, 39--180 and 50--250~MeV, respectively.

The combined study evaluates the actual shared-coupling likelihood over
232 integer masses from 19 to 250~MeV. The active datasets are 2015 at
19--38~MeV, 2015+2016 at 39--49~MeV, all three at 50--90~MeV,
2016+2021 at 91--180~MeV, and 2021 at 181--250~MeV. Independent
constituent random streams are paired into 1,000 joint experiments, with
each whole experiment preserved across the grid. The response basis uses
1,626 support bins and 1,627 deterministic scans. Drawing one joint field
preserves correlations across membership boundaries; combining separate
p-values or treating segments as independent would define another test.

\begin{table}[htbp]\centering\small
\begin{tabular}{lrrrr}\toprule
Profiled scan & Selected mass & GP exceedances & GP global $p$ & Direct exceedances\\
 & [MeV] & out of 200,000 & & out of 1,000\\\midrule
2015 & 51 & 3,797 & 0.018985 & 13\\
2016 & 42 & 0 & $<1.498\times10^{-5}$ & 0\\
2021 10\% & 92 & 5,038 & 0.02519 & 26\\
Combined union & 76 & 0 & $<1.498\times10^{-5}$ & 0\\\bottomrule
\end{tabular}
\caption{Existing conditional significance-GP diagnostics under the
archived coherent stress background, using the principal ordering in
Eq.~\ref{v5-eq:global-ordering}. Zero GP exceedances give the displayed
one-sided 95\% sampling bound; zero direct exceedances give only
$p<0.002991$. These bounds condition on their respective simulation models.
The table reports neither a validated rare particle tail nor a global
calibration of the two-truth envelope in Appendix~\ref{v5-sec:calibration}.}
\label{v5-tab:global-diagnostics}
\end{table}

The latest probability audit gives a particularly instructive combined
example. At 76~MeV the observed $r=0.166$ corresponds to an asymptotic local
$p_0=0.434$. The stress background predicts $a_m=-8.700$ with width 0.979,
so its offset-adjusted score is 9.053. This large score measures a
mismatch with that selected background reference. It cannot be reported
as a $9\sigma$ particle signal. Under the alternative raw ordering, every
simulated scan exceeds the observed maximum and the estimated global
probability is one. Neither comparison is a general goodness-of-fit test.

Direct scans check central distributions, widths, correlations and
accessible maximum tails. In the 2016 profiled study one unadjusted
bulk-distribution comparison has nominal $p=0.034$, reinforcing the need
to inspect approximation discrepancies. Agreement elsewhere and the
absence of adjusted marginal flags do not validate the unresolved rare
tails. The 2016 stress shape also retains source-fit and numerical
qualifications discussed in Appendix~\ref{v5-sec:calibration}.

\enlargethispage{\baselineskip}
\paragraph{Remaining requirements.}
Final particle-global interpretation requires a justified background family
and kernel policy, independent validation of local probabilities, and adequate
simulation of the complete mass and model search. Pending those checks,
Table~\ref{v5-tab:global-diagnostics} reports conditional diagnostics and
sampling bounds.
